% !TeX program = lualatex
% Compile twice from tubs-latex-cd-lite-main (the TU class uses relative asset paths):
% lualatex ../latex/projektarbeit_proposal/v2/schedule.tex
% In the TU Overleaf project, use this file as main.tex and select LuaLaTeX.
% TU template: © 2025 Technische Universität Braunschweig – Stabsstelle für Presse und Kommunikation
\documentclass[basedocumentclass=scrartcl,font=nexus,papersize=a4paper,
    fontsize=11,sideoption=oneside,language=english,tubanner=topleft,
    redDivisionLine=fullpage,template=off]{tubsCD/tubscd}
\setsansfont[Path=./tubsCD/fonts/Nexus/,BoldFont=NexusSansPro-Bold.otf,
    ItalicFont=NexusSansPro-Italic.otf,BoldItalicFont=NexusSansPro-BoldItalic.otf]
    {NexusSansPro-Regular.otf}
\myInstitute{}{}
\usepackage{hyperref}
\hypersetup{hidelinks,pdftitle={Self-Imitation Learning for Robotic Cube Stacking: Project Schedule},
    pdfauthor={Artur Kück},pdflang={en}}
\newcommand{\blockheading}[1]{\par\medskip\noindent{\sffamily\large\bfseries #1}\par\nobreak\smallskip}
\newcommand{\milestone}[3]{\par\medskip\noindent{\sffamily\bfseries #1\quad #2}\hfill{\small\itshape #3}\par\nobreak}
\newcommand{\projectheader}[1]{\noindent{\small\sffamily #1\hfill Artur Kück}\par\medskip}
\begin{document}
\pagestyle{empty}
\setlength{\parindent}{0pt}
\drawSigil{%
\thispagestyle{empty}%
\par\noindent\hfill\begin{minipage}[t]{.88\linewidth}
\projectheader{Self-Imitation Learning for Robotic Cube Stacking}
{\LARGE\sffamily\bfseries Work plan and milestones\par}
\medskip Three months from the project start. I and II denote the first and second halves of each month.
\par\bigskip
\noindent\resizebox{\linewidth}{!}{%
\begin{tikzpicture}[x=1.5cm,y=.51cm,font=\small\sffamily]
\foreach \x/\m in {1/1,3/2,5/3}{\node at (\x,1.85) {Month \m};}
\foreach \x/\halfname in {0.5/I,1.5/II,2.5/I,3.5/II,4.5/I,5.5/II}{\node at (\x,1) {\halfname};}
\foreach \x in {0,...,6}{\draw[black!22] (\x,.55)--(\x,-6.5);}

\node[anchor=east] at (-.15,0) {Literature and protocol}; \fill[tuBlue] (0.04,-0.23) rectangle (0.96,0.23);
\node[anchor=east] at (-.15,-1) {Collect and replay segments}; \fill[tuBlue] (0.04,-1.23) rectangle (1.96,-0.77);
\node[anchor=east] at (-.15,-2) {Train and validate critic}; \fill[tuBlue] (2.04,-2.23) rectangle (2.96,-1.77);
\node[anchor=east] at (-.15,-3) {Train self-imitation policies}; \fill[tuBlue] (2.04,-3.23) rectangle (3.96,-2.77);
\node[anchor=east] at (-.15,-4) {Evaluate on held-out starts}; \fill[tuBlue] (4.04,-4.23) rectangle (4.96,-3.77);
\node[anchor=east] at (-.15,-5) {Write, revise and submit}; \fill[tuBlue] (3.04,-5.23) rectangle (5.96,-4.77);
\node[anchor=east,font=\small\bfseries] at (-.15,-6) {Milestones};
\node at (0.5,-6) {M1};
\node at (1.5,-6) {M2};
\node at (3.5,-6) {M3};
\node at (4.5,-6) {M4};
\node at (5.5,-6) {M5};
\end{tikzpicture}}\par\bigskip
\milestone{M1}{Protocol fixed}{Month 1, midpoint}
Specify the exploration bounds, data splits, success criterion and training budget. Verify that replay reproduces the baseline outcome from the same state.
\milestone{M2}{Dataset frozen}{End of month 1}
Archive the paired segments and record counts of rescues, losses and ties. If rescues are rare, quantify the attainable gain and retain that limitation in the evaluation.
\milestone{M3}{Models frozen}{End of month 2}
Assess critic rankings on separate start families. Save all three BC variants for each training seed before opening the final test set.
\milestone{M4}{Evaluation complete}{Month 3, midpoint}
Report paired success-rate differences against uniform replay and the frozen baseline, with results by training seed.
\milestone{M5}{Project submitted}{End of month 3}
Deliver the report, experiment protocol and scripts needed to reproduce the comparisons.
\end{minipage}\hfill\null\par}
\end{document}
